\subsection{Cel projektu}

\paragraph{}Celem niniejszego projektu jest analiza wpływu częściowej okluzji twarzy, ze szczególnym uwzględnieniem zasłonięcia obszaru oczu, na skuteczność systemów rozpoznawania tożsamości w zadaniu identyfikacji osób. Praca koncentruje się na ocenie odporności współczesnych modeli głębokiego uczenia na zakłócenia wizualne typowe dla rzeczywistych scenariuszy, takich jak monitoring wizyjny czy zastosowania mobilne.

W ramach projektu zaprojektowano i zaimplementowano modyfikacje architektury bazowej IResNet50 (ArcFace), polegające na integracji mechanizmu uwagi typu Convolutional Block Attention Module (CBAM) w blokach rezydualnych oraz na zastosowaniu pomocniczej gałęzi Transformer jako dodatkowej funkcji straty. Celem tych działań było zwiększenie zdolności modeli do selektywnej ekstrakcji cech dyskryminacyjnych w warunkach częściowej utraty informacji wizualnej.

Skuteczność zaproponowanych rozwiązań oceniono poprzez porównanie z wybranymi modelami typu state-of-the-art, w tym zarówno ciężkimi architekturami referencyjnymi, jak i ich lekkimi wariantami przeznaczonymi do zastosowań o ograniczonych zasobach obliczeniowych. Przeprowadzone eksperymenty umożliwiają sformułowanie wniosków dotyczących wpływu okluzji na proces identyfikacji oraz efektywności architektonicznych modyfikacji w kontekście zwiększania odporności modeli na tego typu zakłócenia.